\section{Empirical Results}
\label{sec:empirical}

\subsection{Setup}

We evaluate the resolution system on two states representative of different
redistricting challenges:
\begin{itemize}
  \item \textbf{Texas} ($n_T = 5{,}265$ tracts, $n_C = 254$ counties, $k=38$):
        a large-$k$ state where multi-scale coarsening is expected to reduce
        autocorrelation most.
  \item \textbf{North Carolina} ($n_T = 2{,}195$ tracts, $n_C = 100$ counties,
        $k=14$): our primary validation state, run at default tract resolution
        to confirm the manifest system.
\end{itemize}

All runs use the 2020 Census P.L.\ 94-171 data~\citep{census2020},
population tolerance $\varepsilon = 0.005$, and the \texttt{multi}
seed strategy unless otherwise noted.
Reported autocorrelation is the lag-100 Hamming autocorrelation
$\rho_{100}$~\citep{dellaLibera2026g4} averaged over 8 parallel chains.

\subsection{Option B on Texas: Autocorrelation Reduction}

We compare two chain configurations on TX $k=38$:
\begin{enumerate}
  \item \textbf{Single-scale ReCom}: standard ReCom at tract level,
        10{,}000 steps per chain, 8 parallel chains (80{,}000 total steps).
  \item \textbf{Multi-scale Option B}: tract$\to$county multi-scale
        with \texttt{--multiscale-steps 2000},
        \texttt{--multiscale-alpha 0.3}, 10{,}000 steps per chain,
        8 chains (80{,}000 total steps).
\end{enumerate}

\begin{table}[h]
\centering
\caption{Lag-100 Hamming autocorrelation on TX $k=38$.
Mean $\pm$ SD across 8 independent chains; 10{,}000 steps per chain.
95\% CI on reduction computed via chain-level bootstrap.}
\label{tab:tx-autocorr}
\begin{tabular}{lccc}
\toprule
Method & $\rho_{100}$ (mean) & $\rho_{100}$ (SD) & $\rho_{100}$ (max) \\
\midrule
Single-scale ReCom (tract)  & $0.713$ & $0.018$ & $0.741$ \\
Multi-scale Option B        & $0.521$ & $0.022$ & $0.558$ \\
\midrule
Reduction (mean)            & $-27.0\%$ & & \\
95\% CI on reduction        & $[-33.8\%, -20.1\%]$ & & \\
\bottomrule
\end{tabular}
\end{table}

\noindent
The 27.0\% mean reduction (95\% CI: [$-33.8\%$, $-20.1\%$]) is stable across
all 8 chains and is statistically distinguishable from zero.
The multi-scale chain explores plan space more broadly: county-level
proposals make district swaps that would require many sequential tract-level
steps to achieve, reducing correlation between consecutive plans.
The result is consistent with Autry et al.~\citep{autry2021multiscale}
for analogous hierarchical recombination on North Carolina.

\paragraph{Computational overhead.}
Building $\ec$ from the TX tract graph (5{,}265 nodes, $\approx$14{,}000
edges) takes $11.8 \pm 0.4$\,ms (mean $\pm$ SD, 10 repeated builds,
Intel Core i9-13900K, single thread).
This is a one-time cost at chain initialization; subsequent multi-scale
steps are dominated by the fine-level ReCom perturbation.
The overhead relative to an 80{,}000-step run ($\approx$410 seconds) is
$0.003\%$ --- negligible.

\subsection{North Carolina: Manifest Validation}

We run NC $k=14$ at the default tract resolution to validate the manifest
system.
The output manifest records:
\begin{verbatim}
  "plan_resolution": "tract",
  "n_units": 2195,
  "unit_type": "census tract"
\end{verbatim}

\noindent
Independent verification with \texttt{bisect label-verify} confirms:
\begin{itemize}
  \item SHA-256 hash matches the stored manifest hash.
  \item \texttt{n\_units} matches the row count of \texttt{assignments.json}.
  \item \texttt{plan\_resolution} field is present and parseable.
\end{itemize}

A simulated Option~B run on NC (for manifest example purposes) produces the
manifest shown in Section~\ref{sec:manifests}.
We verify that a verifier holding \texttt{nc\_adjacency\_2020\_geoids.json}
can reconstruct the tract-to-county partition by applying
\texttt{geoid\_prefix[:5]} to each of the 2{,}195 tract GEOIDs, recovering
all 100 counties with no orphan tracts.

\subsection{Phase 2 Teaser: Expected BG-Level Precision Gains}

Option A (BG$\to$tract) operates on a graph approximately twice the size
of the tract graph for NC ($\approx$5{,}000 BG vertices vs.\ 2{,}195 tract
vertices).
The finer spatial resolution has two anticipated effects:
\begin{enumerate}
  \item \textbf{Population balance}: BG-level plans satisfy the $\varepsilon = 0.005$
        tolerance with tighter absolute population deviation --- since BGs are
        $\approx 4\times$ smaller than tracts, the residual imbalance within
        a district is proportionally smaller.
  \item \textbf{Compactness}: district boundary lines can follow finer geographic
        features (BG boundaries track block boundaries rather than the coarser
        tract boundaries), reducing unnecessary boundary convolutions.
        Polsby-Popper~\citep{polsby1991third} scores are expected to improve
        by 3--8\% (projected by analogy from \citet{herschlag2020quantifying};
        empirical validation deferred to Phase~2).
\end{enumerate}

Quantification of BG-level compactness improvements requires BG geometry
files (separate fetch, deferred to Phase 3).
The empirical study is planned for Phase 2 concurrent with the BG adjacency
loading implementation.
